\scanpage{33}
\begin{center}
1977/78 Diploma of Advanced Studies course by A. Grothendieck

\uline{Arithmetic geometry of regular polygons}
\end{center}

{\reversemarginpar\margin{\handwritten{!}}\normalmarginpar}%
% The hand of the marginal ! and short underline is unassigned.
% The short hand-drawn rule underlines the typed words en commun.
This ``course'' will take place in the spirit of a \handwritten{\uline}{joint} investigation of multiple facets of the notion of a regular polygon over a base field, then over an arbitrary base ring. This will be a stimulating opportunity to familiarize ourselves with the basic notions of vector, \tikz[remember picture,overlay]{\coordinate (s033-affine);}affine, projective geometry - with or without a metric - in dimensions one, two and three, and with the phenomena particular to these dimensions, as well as with the sometimes unexpected aspects taken on by familiar notions (such as that \tikz[remember picture,overlay]{\coordinate (s033-de);}of reflection\deleted{s}, or of the centre of a regular polygon) over fields of positive characteristic, and in particular over finite fields. The need for a geometry over general base rings (such as the ring ``over which is written'' ``the most general'' regular $\mathrm{n}$-polygon) will make itself felt along the way, and will probably lead us to touch on a few\\
\tikz[remember picture,overlay]{\coordinate (s033-points1);}points in the foundations of algebraic geometry. If time permits and if other investigations do not detain us , we shall attempt, on the fly, to construct the ``general'' regular icosahedron over a suitable base ring, finite over the ring \hspace{.5em}${\handwritten{\mathbb}{Z}}$\hspace{.5em} of integers, having residue fields of all characteristics. (Characteristic two \hspace{.5em}will probably give us a hard time !).

From an intuitive and familiar starting point, this ``course'' would like\tikz[remember picture,overlay]{\coordinate (s033-voudrait);} to be an opening onto a fascinating universe, full of mysteries which ask only to be elucidated. It can bear fruit only with \tikz[remember picture,overlay]{\coordinate (s033-active);}active participation by the student and substantial personal work, in particular to develop w\deleted{e}ith\\
\tikz[remember picture,overlay]{\coordinate (s033-detail);}the necessary detail the many points\handwritten{(}of foundation\inserted{\handwritten{s}} in particular) which will only have been sketched in the ``course''. He is very strongly advised to attend at the same time the course \hspace{.5em}C4\hspace{.5em} on the icosahedron, which will be held in a similar spirit. There will probably be numerous interferences between these two courses of reflection.

% The underlying e in avec has a visible mid-height horizontal bar.
% The live v is clear; the heavy overstrike leaves its inscription method open.
% The opening parenthesis retains C's hand encoding; reinforcement of type remains open.
% Z has two obliques; the added diagonal retains hand encoding at medium confidence,
% with a second type impression remaining an instrument alternative.
% C4 characters are clear; the tall narrow C and descending unfooted 4 leave
% inscription method unresolved. Plain encoding is retained without a glyph doubt.

% Detached geometries: no verbal role, origin or author is assigned.
% Local source-pixel geometry is scaled to the typeset x-height; the zero-width
% text anchors retain each mark's relation to the corresponding passage.
% The line break before le detail retains the left-edge mark's line attachment.
\figure{\begin{tikzpicture}[remember picture,overlay]
% Pale oval/U-shaped trace below the a of affine, full-source area x=1255--1264, y=638--650.
\begin{scope}[shift={(s033-affine)},x=0.26pt,y=-0.26pt]
\draw[black!30,line width=0.65pt,line cap=round] (3,0) .. controls (1,4) and (2,9) .. (6,10) .. controls (9,9) and (10,5) .. (8,2);
\end{scope}
% Compact point and paler broken trace below de before réflexion, source (1019,877).
\begin{scope}[shift={(s033-de)},x=0.26pt,y=-0.26pt]
\fill[black!45] (19,19) ellipse[x radius=1.25,y radius=.95];
\draw[black!25,line width=.45pt] (25.5,21.7) .. controls (27,20.9) and (29.5,21.3) .. (30.7,22.7);
\end{scope}
% Pale taper beside the source line beginning points, source (55--60,1220--1224).
\begin{scope}[shift={(s033-points1)},x=0.26pt,y=-0.26pt]
\fill[black!40] (-199,7) -- (-194.5,9.6) -- (-198.8,9.9) -- cycle;
\end{scope}
% Detached dark hook below the end of voudrait, around x=1498, y=1582.
\begin{scope}[shift={(s033-voudrait)},x=0.26pt,y=-0.26pt]
\draw[line width=0.75pt,line cap=round] (-17,16) .. controls (-15,20) and (-19,22) .. (-21,22);
\end{scope}
% Detached upper-right / lower-left pair below active, around (1534,1722), (1529,1729).
\begin{scope}[shift={(s033-active)},x=0.26pt,y=-0.26pt]
\fill (31,16) ellipse[x radius=1.5,y radius=1.1,rotate=45];
\fill (26,23) ellipse[x radius=1.3,y radius=1.0,rotate=20];
\end{scope}
% Compact two-lobed mark at the physical left edge beside le detail, around x=28--37,y=1839--1846.
\begin{scope}[shift={(s033-detail)},x=0.26pt,y=-0.26pt]
\fill (-227,-5) .. controls (-227,-8) and (-223,-8) .. (-222,-5)
  .. controls (-219,-6) and (-217,-5) .. (-217,-3)
  .. controls (-217,0) and (-220,1) .. (-221,0)
  .. controls (-223,0) and (-223,-3) .. (-224,-3)
  .. controls (-226,-3) and (-228,-3) .. (-227,-5) -- cycle;
\end{scope}
\end{tikzpicture}}
% archive sheet number: 32
